% ============================================================
%  CHAPTER 7 — Discussion
% ============================================================

\chapter{Discussion}
\label{ch:discussion}

% ============================================================
\section{Validity of Results}
% ============================================================

\subsection{Embedded Subsystem}

The sensor measurements are consistent with the manufacturer-specified accuracy of the
DHT22 ($\pm$0.5\,°C, $\pm$2--5\,\% RH) and within the expected operational range of
the MAX30102 for resting-state adult subjects.
The BPM standard deviation of 2.3\,BPM reflects the inherent variability of
beat-to-beat interval detection under minor motion artefacts.
The SpO\textsubscript{2} accuracy of $\pm$1\,\% is within the $\pm$2\,\% tolerance
generally accepted.

\subsection{Backend and AI Subsystems}
The backend procedures are structurally consistent with the tRPC schema and are
guarded by session-based authentication. Quantitative latency and accuracy metrics
for AI inference remain to be measured in the deployment environment. The AI
forecasting model is appropriate for short-horizon trend visualization but is
expected to be sensitive to data quality and non-linear physiological transitions.

\subsection{Web Dashboard Subsystem}
The dashboard operates with real-time updates and a 600\,ms face recognition loop.
Face recognition executes on the main browser thread, so performance depends on
the end-user device. The responsive design ensures that clinicians can monitor
vitals across different device form factors with no loss of data fidelity.

% ============================================================
\section{Limitations}
% ============================================================

\subsection{Embedded Subsystem}

\begin{description}
	\item[No local offline buffer.]
	      Records generated during network outages are discarded.

	\item[Plaintext NVS storage.]
	      Firebase credentials (API key, email, password) are stored in NVS without
	      encryption.
	      The ESP32 NVS partition supports native encryption via an NVS key partition;
	      this should be enabled before any production deployment.

	\item[Implicit shared-memory atomicity.]
	      The sensor globals are 32-bit floats written on Core~1 and read on Core~0
	      without explicit \mbox{synchronisation}.
	      A FreeRTOS mutex or atomic qualifier is the correct long-term solution.

	\item[Single patient per device.]
	      The current design maps one device to exactly one patient and one room.
	      Multi-patient rooms would require either multiple devices or a structural change
	      to the data model.
\end{description}

\subsection{Backend Subsystem}
\begin{description}
	\item[No rate limiting.] The current tRPC implementation does not enforce rate-limiting, which could make the system vulnerable to denial-of-service attacks if exposed to the public internet.
	\item[SMTP delivery latency.] Email alerts are dependent on external SMTP relay availability, which may introduce latencies of 2--5 seconds during network congestion.
	\item[Single-region deployment.] The backend and Firebase services are currently deployed in a single geographic region, increasing latency for cross-regional users and lacking multi-region failover.
\end{description}

\subsection{Web Dashboard Subsystem}
\begin{description}
	\item[High resource consumption.] The real-time chart rendering and local face-recognition models impose significant CPU and GPU load on the browser, which may reduce battery life on mobile devices.
	\item[Browser-specific optimizations.] While cross-browser compatible, the system performs best in Chromium-based browsers due to advanced WebGL optimizations.
	\item[No persistent offline mode.] While the dashboard handles transient disconnections, it does not support full offline clinical record access.
\end{description}

\subsection{AI Subsystem}
\begin{description}
	\item[Linear trend assumptions.] The current forecasting model uses linear regression, which may fail to capture complex non-linear physiological events such as sudden cardiac arrhythmias.
	\item[Short baseline window.] The analysis window is short (target 20 points, minimum 5), providing only local context and limiting robustness during long-term drift.
	\item[Data quality dependency.] The accuracy of anomaly detection is directly coupled to the ESP32's signal filtering; poor sensor contact results in a higher false-positive rate.
\end{description}

% ============================================================
\section{ESP32 Constraints Impact}
% ============================================================

The most significant platform constraint encountered was the non-thread-safety of
the mbedTLS context, which imposed the non-obvious architectural requirement of
deferring Firebase initialisation until the task is running on Core~0.

The 8\,KiB task stack was found to be adequate for the TLS handshake and async send
operations, with an estimated 1.5\,KiB headroom.

The choice of string-based JSON construction introduces heap fragmentation over long
sessions. This risk is partially mitigated by the 1\,Hz push interval but remains a
concern for multi-day deployments.

% ============================================================
\section{Robustness}
% ============================================================

The unified \texttt{wifiConnect()} function using \texttt{WiFi.begin()} handles both
cold-start failure and mid-session dropout correctly, as confirmed by the
\textit{mid-session Wi-Fi dropout} and \textit{cold-start, Wi-Fi delayed} scenarios
reported in Table~\ref{tab:stability}.

The delta-threshold filter provides robustness against spurious sensor noise:
fluctuations smaller than 0.3\,°C, 0.5\,\% RH, or 2\,\% SpO\textsubscript{2} do
not generate database writes, reducing both network load and the volume of
low-information records in the database.

% ============================================================
\section{Identified Improvements}
% ============================================================

\subsection*{Embedded Subsystem}

Priority improvements for future versions of the embedded firmware:

\begin{enumerate}
	\item \textbf{Local offline buffer.}
	      A LittleFS-backed FIFO queue for records generated during network outages,
	      with automatic retry on reconnection.
	      This would eliminate the record loss observed in the dropout and cold-start
	      stability scenarios.

	\item \textbf{NVS encryption.}
	      Enable the ESP32 native NVS encryption partition to protect credentials at rest.

	\item \textbf{Thread-safe sensor globals.}
	      Replace raw float globals with a mutex-protected struct or a FreeRTOS queue
	      for safe inter-core communication.

	\item \textbf{Wi-Fi modem sleep.}
	      Implement modem sleep between upload intervals to reduce average current usage for future battery-powered deployments.
\end{enumerate}

\subsection*{Backend, Web Dashboard, and AI Subsystems}

Priority improvements for the server-side and client-side subsystems:

\begin{enumerate}
	\item \textbf{Deep Learning Models.} Migrating the AI layer to an LSTM (Long Short-Term Memory) or Transformer architecture to better capture non-linear physiological trends and improve forecasting accuracy during acute events.

	\item \textbf{PWA Support.} Implementing Progressive Web App (PWA) features to enable offline data caching and native push notifications for critical clinical alerts, reducing dependency on an always-open browser tab.

	\item \textbf{Distributed Cache.} Adding a Redis-based caching layer for the tRPC backend to accelerate historical data retrieval for long-term patient records and reduce Firestore read costs.

	\item \textbf{Multi-Factor Authentication.} Enhancing security for clinical staff logins using hardware security keys or biometric verification to meet healthcare-grade access control standards.
\end{enumerate}
